% !TEX root = ../main.tex
% Arquitectura de la aplicacion web (IT-42), contada como EL RECORRIDO DE UNA
% PREGUNTA.
%
% Antes era un grafo de componentes con las dependencias colgando a la derecha
% y no se entendia: habia que seguir seis flechas para saber por donde entra
% una pregunta. Aqui no se siguen flechas, se siguen numeros.
%
% Lo que la figura ensena y las anteriores escondian: al modelo se le llama
% DOS veces por turno, una para fijar de que titulacion se habla y otra para
% redactar. Es la mitad de lo que tarda una respuesta. Por eso son esos dos
% pasos, y solo esos dos, los que llevan el borde grueso.
%
% La salida del paso 3 no es un adorno: cuando la pregunta se contesta con
% texto fijo NO se llama al modelo ni una vez. Es la regla de que lo que el
% sistema no puede permitir no se le pide al modelo, se impide antes.
%
% Los pasos van en serpentina ---la fila de abajo se lee de derecha a
% izquierda--- para que las ocho cajas quepan en dos filas y no en una tira de
% ocho, que obligaria a encoger la letra a la mitad.
%
% Las cajas se encadenan con `right=of`, `left=of` y `below=of`: dar las
% posiciones a mano ya provoco solapes en otras tres figuras, porque una caja
% crece con su texto y `minimum height` solo fija el minimo.
\begin{figure}[H]
  \centering
  \resizebox{\textwidth}{!}{%
    \begin{tikzpicture}[
      paso/.style={nodoDiag, text width=3.2cm, minimum height=2.1cm},
      pasoFoco/.style={paso, draw=diagFocoBorde, fill=diagFocoFondo, very thick},
      salida/.style={nodoDiag, fill=white, densely dashed,
        text width=3.2cm, minimum height=1.5cm},
      % El número va CENTRADO en la esquina, medio dentro y medio fuera: puesto
      % dentro se sienta encima de la primera línea de texto.
      num/.style={circle, draw=diagNodoBorde, fill=white, font=\small,
        inner sep=1pt, minimum size=6mm},
      numFoco/.style={num, draw=diagFocoBorde, very thick},
      rotulo/.style={font=\scriptsize, inner sep=2pt},
      ]

      % ── Fila de arriba: de la pregunta a la decisión de ámbito ─────────────
      \node[paso] (p1)
      {El estudiante\\escribe\\{\scriptsize en el navegador}};
      \node[paso, right=0.6cm of p1] (p2)
      {El servidor recibe la pregunta\\{\scriptsize \texttt{/api/chat}}};
      \node[paso, right=0.6cm of p2] (p3)
      {¿Se contesta con texto fijo?\\{\scriptsize saludo · otro centro}};
      \node[pasoFoco, right=0.6cm of p3] (p4)
      {Se decide de qué titulación se habla\\{\scriptsize 1.\textsuperscript{a} llamada al modelo}};

      % ── Lo que corta el recorrido antes de llegar al modelo ────────────────
      \node[salida, below=0.9cm of p3] (fija)
      {Respuesta fija\\{\scriptsize sin llamar al modelo}};

      % ── Fila de abajo: se lee de derecha a izquierda ───────────────────────
      \node[pasoFoco, below=0.9cm of fija] (p6)
      {Se redacta la respuesta\\{\scriptsize 2.\textsuperscript{a} llamada al modelo}};
      \node[paso, right=0.6cm of p6] (p5)
      {Se buscan los fragmentos\\{\scriptsize los que pasan el suelo}};
      \node[paso, left=0.6cm of p6] (p7)
      {Se comprueba lo redactado\\{\scriptsize ninguna titulación}\\{\scriptsize inventada}};
      \node[paso, left=0.6cm of p7] (p8)
      {Llega por partes al navegador\\{\scriptsize y se anota el turno}};

      % ── El recorrido ───────────────────────────────────────────────────────
      \draw[flechaDiag] (p1) -- (p2);
      \draw[flechaDiag] (p2) -- (p3);
      \draw[flechaDiag] (p3) -- node[rotulo, above=2pt]{no} (p4);
      \draw[flechaDiag] (p3) -- node[rotulo, right=2pt]{sí} (fija);
      \draw[flechaDiag] (p4) -- (p5);
      \draw[flechaDiag] (p5) -- (p6);
      \draw[flechaDiag] (p6) -- (p7);
      \draw[flechaDiag] (p7) -- (p8);

      % ── Los números, después de las cajas para quedar por encima ───────────
      \node[num]     at (p1.north west) {1};
      \node[num]     at (p2.north west) {2};
      \node[num]     at (p3.north west) {3};
      \node[numFoco] at (p4.north west) {4};
      \node[num]     at (p5.north west) {5};
      \node[numFoco] at (p6.north west) {6};
      \node[num]     at (p7.north west) {7};
      \node[num]     at (p8.north west) {8};

      % ── Dónde acaba una fase y empieza la otra ─────────────────────────────
      % La nota se centra bajo la fila entera, no bajo una caja: el punto medio
      % entre los dos extremos sale de un `midway` sobre un camino invisible.
      \path (p8.south) -- (p5.south) node[midway] (abajo) {};
      \node[rotulo, anchor=north, text width=13cm, align=center]
      at ([yshift=-0.9cm]abajo.center)
      {Los pasos 1 a 3 y el 8 son la aplicación web; del 4 al 7 es el flujo de
        la Fase 2, que se importa como biblioteca y no se modifica. Una
        respuesta cuesta dos llamadas al modelo; cuando el paso 3 corta, no
        cuesta ninguna.};
    \end{tikzpicture}%
  }
  \caption[Arquitectura de la aplicación web]{\textcolor{borrador2}{El recorrido
      de una pregunta por la aplicación. Se sigue por los números y no por las
      flechas. Los dos pasos con el borde grueso son los que llaman al modelo
      generativo: son dos por turno, uno para fijar de qué titulación se habla
      y otro para redactar. La salida del paso~3 es la que evita los dos: una
      pregunta que se contesta con texto fijo no llega al modelo. Fuente:
      Elaboración propia.}}
  \label{fig:arquitectura_web}
\end{figure}
